\section{Discussion}
\label{sec:discussion}

\subsection{Interpretation of Detected Sequences}
\label{subsec:sequences_found}

Where the Monte~Carlo permutation test returns $p < 0.01$ after the
Benjamini--Hochberg false discovery rate (FDR) correction
\citep{BenjaminiHochberg1995}, the identified sequences represent
statistically robust departures from a spatially inhomogeneous Poisson
null model.  Three physical mechanisms can account for inter-event
coupling over the tectonic distances recovered here.

\paragraph{(a) Static Coulomb stress transfer.}
Changes in static Coulomb failure stress ($\Delta\mathrm{CFS} \gtrsim
0.01$~MPa) generated by a large rupture can advance or retard failure on
neighbouring fault segments within $\sim\!100$~km of the source
\citep{King1994, Stein1999}.  The effect is permanent (quasi-static)
and decays with distance as $r^{-3}$ in a homogeneous elastic half-space.
Triggered sequences consistent with this mechanism are therefore expected
to cluster within years of the mainshock and to respect tectonic
connectivity, a signature directly targeted by our $\eta$-metric.  The
rate-and-state frictional framework of \citet{Dieterich1994} predicts
that seismicity rate enhancements persist for a nucleation time
$t_a \propto \dot{\tau}^{-1}$, typically months to $\sim\!10$~years for
interplate environments, broadly matching the 1–5-year temporal window in
which detected sequences are most stable (see Section~\ref{sec:robustness}).

\paragraph{(b) Dynamic stress transfer.}
Transient dynamic stresses carried by surface waves ($\Delta\sigma_\mathrm{dyn}
\sim 10^{-3}$–$10^{-1}$~MPa) can trigger seismicity at distances of
thousands of kilometres within seconds to hours of the causative rupture
\citep{Hill1993, Gomberg2001, Brodsky2006}.  The attenuation of dynamic
stresses is slower than static stresses ($r^{-1}$ to $r^{-2}$), making
this mechanism the dominant candidate for ultra-distal coupling across
plate boundaries.  However, because dynamic triggering typically manifests
on timescales far shorter than our minimum resolution ($\Delta t_{\min}
\sim 1$~day), sequences whose inter-event intervals fall in the hours-to-days
range may have a dynamic origin.  Disentangling this contribution requires
waveform-level analysis beyond the scope of the present catalogue study.

\paragraph{(c) Viscoelastic postseismic relaxation.}
Viscous flow of the lower crust and asthenosphere following large
ruptures redistributes stress over timescales of months to decades
\citep{Pollitz1998, Freed2001, Freed2005}.  Unlike static stress
changes—which are immediate—viscoelastic loading grows with time before
eventually decaying, potentially elevating failure probabilities for
subsequent events years after the mainshock.  This mechanism is
particularly relevant for events on oceanic plates where asthenospheric
viscosity is low ($\eta \sim 10^{18}$–$10^{19}$~Pa\,s;
\citealt{Pollitz1998}).  Sequences with inter-event intervals of 1–10
years that span multiple tectonic zones may thus reflect postseismic
rather than co-seismic loading.

Cascading effects described by \citet{MarsanLengliene2008} further
complicate attribution: intermediate $M\geq 6.5$ events that are
themselves triggered may relay stress to regions beyond the direct
influence of the initial mainshock, effectively extending the network of
influence along plate-boundary paths.  The ETAS-based analysis proposed
in Section~\ref{subsec:perspectives} is the natural next step for
distinguishing direct from relay-mediated coupling
\citep{HelmstettterSornette2003, Ogata1998}.

\subsection{Interpretation in the Absence of Significant Sequences}
\label{subsec:sequences_absent}

A null result—no sequences surviving FDR correction at $q = 0.05$—does
not imply that global seismicity is Poissonian.  It implies, more
precisely, that our method, applied to the available catalogue with the
chosen thresholds ($N \geq 4$ events, $\geq 3$ tectonic zones, $T_w \leq
5$~years), lacks the sensitivity to detect departures from the null at
this significance level.

Three sources of negative bias deserve explicit acknowledgement.
\emph{Catalogue incompleteness} prior to $\sim\!1900$~CE inflates
inter-event intervals and dilutes the density of the nearest-neighbour
distance distribution, reducing the power of the $\eta$-metric to
separate clustered from background pairs.  The $\mathrm{quality\_score}$
weighting partially mitigates this by down-weighting pre-instrumental
events, but cannot compensate for missing events entirely.
\emph{Temporal non-stationarity}, arising from the exponential growth in
the number of well-located events since the establishment of global
seismic networks in the 1960s, inflates the apparent event rate in the
modern period and may cause the permutation test to underestimate
$p$-values for recent sequences relative to historical ones.  Finally,
\emph{suboptimal temporal windows}: the characteristic timescale of
inter-plate stress transfer is poorly constrained and may differ from the
windows explored here; an exhaustive sensitivity analysis (Section~\ref{sec:robustness})
is therefore essential before a definitive null result can be stated.

\subsection{Methodological Limitations and Their Mitigation}
\label{subsec:limitations}

We identify six principal limitations and describe the mitigation strategy
adopted for each.

\begin{enumerate}

\item \textbf{Threshold arbitrariness ($N \geq 4$, $\geq 3$ zones).}
The minimum sequence length and zone-span requirements were chosen to
limit false positives but have no strong physical motivation.  We address
this through synthetic catalogue experiments (Section~\ref{sec:robustness},
Experiment~3.5) in which known sequences of varying length are injected;
thresholds yielding $F_1 > 0.7$ are adopted as empirically supported.

\item \textbf{Absence of multiple-testing correction in naive analysis.}
Testing $K$ candidate sequences simultaneously inflates the family-wise
error rate.  We apply the linear step-up procedure of
\citet{BenjaminiHochberg1995} at $q = 0.05$, which controls the expected
proportion of false discoveries rather than the probability of any false
discovery, thereby preserving statistical power relative to Bonferroni
correction when $K$ is large.

\item \textbf{Historical catalogue incompleteness.}
Events before $\sim\!1900$ lack instrumental waveform constraints; their
magnitudes and locations carry larger uncertainties.
\citet{Wiemer2000} and \citet{Woessner2005} provide magnitude-of-completeness
($M_c$) estimates as a function of epoch; we stratify the catalogue at
$M_c$ thresholds for each 50-year bin and restrict global sequence tests
to events above the applicable $M_c$, ensuring that the null model is
conditioned on the same completeness level.

\item \textbf{Non-stationarity of the recording process.}
The raw event rate increases monotonically with time due to network
improvements, not geophysical acceleration.  We normalise observed event
counts to an expected intensity derived from a time-varying NHPP
(non-homogeneous Poisson process) fitted to the decadal event rate, so
that Monte~Carlo permutations preserve the temporal intensity profile of
the real catalogue \citep{Michael2011, ShearerStark2012}.

\item \textbf{No declustering applied to baseline catalogue.}
Retaining aftershocks inflates sequence counts.  In Experiment~3.3 we
compare results with and without declustering via (i)~the
\citet{Gardner1974} window method and (ii)~stochastic declustering using
the ETAS-based approach of \citet{Zhuang2002}.  Consistency across
these variants would indicate that the detected global sequences are not
artefacts of local aftershock contamination.

\item \textbf{Focal depth and mechanism not incorporated.}
Strike-slip, thrust, and normal faulting produce markedly different
$\Delta\mathrm{CFS}$ lobes; shallow crustal events couple differently
from deep subduction events.  The current metric treats all $M \geq 6.5$
events identically.  Incorporating focal mechanism solutions and depth
constraints—available for $\sim\!70\%$ of the instrumental-period
catalogue \citep{Engdahl1998}—is identified as the primary direction for
future methodological development.

\end{enumerate}

\subsection{Comparison with Previous Studies}
\label{subsec:comparison}

Our findings must be situated within the debate initiated by
\citet{Michael2011}, who demonstrated that apparent clustering of $M \geq
7$ events in 1995–2011 was statistically indistinguishable from random
Poisson fluctuations once the inhomogeneous event rate was properly
conditioned.  \citet{ShearerStark2012} independently confirmed that
global $M \geq 7$ and $M \geq 8$ rates showed no significant temporal
increase since the 2004 Sumatra earthquake, ruling out a secular trend in
global hazard.

Our approach addresses both critiques on methodological grounds.  First,
we condition the null model on the empirical spatial intensity of the
catalogue (via tectonic zone stratification), mitigating the rate
inhomogeneity problem raised by \citet{Michael2011}.  Second, we do not
test for a secular trend in global activity—a hypothesis we share with
\citet{ShearerStark2012}—but rather for episodic multi-event sequences
localised in both tectonic space and time, a fundamentally different null
hypothesis.

The landmark observation of \citet{Hill1993} following the 1992 Landers
earthquake established that $M \geq 7$ events can trigger seismicity at
distances exceeding 1\,000~km within hours.  Subsequent work by
\citet{Brodsky2006}, \citet{Gomberg2001}, and \citet{Parsons2002}
extended this phenomenology across multiple tectonic environments.  Our
study asks whether the cumulative record of such interactions over
4\,000~years leaves a detectable imprint in the global $M \geq 6.5$
catalogue—a question that requires the long temporal baseline provided
by the integrated USGS~ComCat / ISC~Bulletin / NOAA~NGDC catalogue and
the tectonic-distance metric we introduce.

The doublet analysis of \citet{KaganJackson1999}, which demonstrated
elevated probabilities of a second large event near the epicentre of a
first within several years, provides perhaps the closest precedent for
our approach.  We extend their framework from pairs to $N \geq 4$
sequences and from Euclidean to tectonic distance, yielding a more
sensitive detection scheme while explicitly accounting for plate-boundary
geometry \citep{Bird2003}.

\subsection{Perspectives and Future Work}
\label{subsec:perspectives}

The most immediate extension of this work is formal ETAS validation.
Fitting a space-time ETAS model \citep{Ogata1998, HelmstettterSornette2003}
to the declustered catalogue and comparing its likelihood with that of
the null Poisson model would provide a model-selection test
(e.g.,~AIC or BIC) that is independent of the permutation framework
developed here.  Sequences that survive both tests would constitute
particularly strong evidence for non-Poissonian global coupling.

The second priority is causal inference through $\Delta\mathrm{CFS}$
modelling.  For each candidate sequence, computing the Coulomb failure
stress imparted by all preceding events in the sequence on the subsequent
rupture plane \citep{King1994, Toda2005} would allow us to distinguish
physically plausible trigger–receiver geometries from geometrically
random co-occurrences.  This requires focal mechanism solutions and fault
geometry models for the majority of events, a task that is feasible for
post-1976 events using the Global CMT catalogue but challenging for the
pre-instrumental period.

Finally, the methodology developed here is directly applicable to
paleoseismic records \citep{McCalpin2009, Griffin2020, Giri2026}, where
individual fault-segment rupture histories extending over $10^3$–$10^4$
years provide the long temporal baselines needed to characterise the
tail of the inter-event time distribution and to test whether global
connectivity persists on millennial timescales.
